\section{Introduction et guide de lecture}\label{int:intro}

\subsection{La question}
Former un professionnel de santé prend de trois à dix ans. Une décision prise aujourd'hui sur les places de
formation, l'âge de départ, l'accueil des diplômés étrangers ou la régulation de l'installation n'agit donc sur
l'offre de soins qu'à l'horizon d'une décennie, et parfois dans d'autres territoires que ceux visés. Les projections
publiées (DREES) sont nationales ou régionales et tendancielles. Elles ne permettent pas de composer ses propres
scénarios.

Le projet construit un outil de recherche qui répond à la question suivante : \emph{quel effectif d'une profession,
dans un territoire donné, à une date donnée, sous une politique donnée ?} La réponse prend la forme d'une
fourchette, publiée avec sa couverture mesurée (aucune bande n'est calibrée à son niveau nominal), et non d'un
chiffre unique. Elle est comparée à la tendance, et l'écart est décomposé par flux.

\subsection{Périmètre et définitions}\label{int:perimetre}
\begin{itemize}
\item \textbf{Professions.} Sages-femmes (SF), chirurgiens-dentistes (CD), pharmaciens (PH),
  masseurs-kinésithérapeutes (MK) et infirmiers (IDE). La médecine générale est exclue par le sujet. Les médecins
  spécialistes (cardiologie, dermatologie, ophtalmologie) sont prévus en phase~2 et ne sont pas modélisés ; le
  jumeau contient seulement une série « spécialistes » de test héritée de la v1, sans validation.
\item \textbf{« Libéral ».} Dans tout ce document, le mot désigne le \emph{secteur privé}, lucratif ou non. Sont
  inclus :
  \begin{itemize}
  \item les libéraux et les mixtes ;
  \item les adjoints salariés d'officine ;
  \item les salariés des établissements privés.
  \end{itemize}
  Tout employeur public est exclu. Les sources emploient « libéral » au sens d'exercice indépendant ; l'écart est
  traité au chapitre~\ref{sec:donnees-ref}.
\item \textbf{Échelles.} Commune, EPCI, bassin de vie, département, région et France. Le calcul descend à la
  commune, mais l'affichage et la validation s'arrêtent au département (chapitre~\ref{sec:validation-ref}).
\item \textbf{Horizon.} 2050 est l'horizon de référence des scénarios. L'horizon jugé par backtest est de
  5~ans ; à 8 et 10~ans, la persistance fait mieux qu'A2\mesure{} (chapitre~\ref{sec:menaces-ref}). Au-delà, les
  sorties sont des \emph{scénarios conditionnels}, pas des prévisions.
\item \textbf{Données.} Uniquement des données ouvertes pour les tables traitées. Chaque fichier traité porte sa
  source, son millésime et le script qui le produit (manifeste contrôlé par \texttt{v2/data/manifest.py --check}).
  Le code exécuté (jumeau, carte, assistant, sensibilité) ne lit lui aussi que des données ouvertes : depuis le
  15 novembre 2026, le vivier des dentistes et des pharmaciens vient de SISE (MESR) et non plus d'un fichier
  d'inscrits fourni à l'équipe sous accord de confidentialité ; le contrôle du manifeste couvre ce code
  (chapitre~\ref{sec:repro-ref}).
\end{itemize}

\subsection{L'état en un coup d'œil}\label{int:etat}
\begin{center}\small
\begin{tabular}{@{}p{0.27\textwidth}p{0.66\textwidth}@{}}
\toprule
Composant & État au 15 novembre 2026 \\
\midrule
Jumeau (chap.~\ref{sec:jumeau-ref}) & 18 critères d'acceptation sur 18\mesure{} (composition A0 ; C5, C9 et la
  variante filière exigent un fichier non commité) ; trois revues adversariales, qui n'ont couvert que C1 à C13 :
  C14 à C18 n'ont été soumis à aucune revue. Un \texttt{run()} direct d'A0 au département prend
  16\,ms\mesure{} (7 séries : les 5 professions, les infirmiers scindés en libéraux et salariés, et une série
  « spécialistes » de test héritée de la v1) ; un scénario à tous leviers servi avec C1 prend 280 à
  330\,ms\mesure. \\
Méthodes (chap.~\ref{sec:methodes-ref}) & C1 (cohortes) servi par la carte et l'assistant, par injection de ses
  taux dans le jumeau : MAPE départementale à $h=5$ de 7,76\,\%\mesure{} [7,12 ; 8,40]. Budget d'erreur : plancher
  de bruit de comptage de 3,11 points\derive{}, borne basse de la part irréductible. Intervalles recalibrés,
  GBM et MLP rejetés. \\
Validation (chap.~\ref{sec:validation-ref}) & 12 protocoles pré-enregistrés. Le grain sous le département, les
  leviers et 2050 ne sont pas validés. Dépôt de prévisions 2027 à juger en 2027. \\
Données et paramètres & 34 fichiers traités commités sous manifeste. Hyperparamètres : 29 fiches relues sur le
  texte intégral des sources le 15 novembre (chap.~\ref{sec:parametres-ref}). \\
Interface & Carte (9 leviers de politique et 2 hypothèses de contexte, pris parmi les 12 paramètres-leviers du
  jumeau), phrase en langage naturel compilée de façon
  déterministe, assistant LLM ancré sur les sorties. \\
Conformité & Pas un dispositif médical ; AIPD probablement non requise ; un point réel sur les petits effectifs
  publiés (chap.~\ref{sec:conformite-ref}). \\
Module de besoin & Programme défini, non construit (chap.~\ref{sec:ouvertes-ref}). \\
\bottomrule
\end{tabular}
\end{center}

\subsection{Ce qui peut être cité, et ce qui ne doit pas l'être}\label{int:citer}
\begin{itemize}
\item \textbf{Peut être cité.}
  \begin{itemize}
  \item Les erreurs de backtest au département, horizon 1 à 5 ans, avec leur intervalle. Ces intervalles sont
    probablement trop étroits : les origines ne sont pas rééchantillonnées (chapitre~\ref{sec:menaces-ref}).
  \item Les écarts entre scénarios, présentés comme des effets \emph{conditionnels au modèle}.
  \item Les comptages et séparations public/privé de l'année de base, avec leur fourchette.
  \end{itemize}
\item \textbf{Ne doit pas être cité comme prévision.}
  \begin{itemize}
  \item Un effectif en 2050, y compris dans les supports de présentation du projet. La reproduction des
    projections DREES à long terme échoue (protocole N2, chapitre~\ref{sec:validation-ref}).
  \item Un effectif communal.
  \item L'effet d'un levier présenté comme causal et mesuré.
  \end{itemize}
\end{itemize}

\subsection{Guide de lecture selon le lecteur}
\begin{itemize}
\item \textbf{Statisticien.} Méthodes (\ref{sec:methodes-ref}), puis validation (\ref{sec:validation-ref}),
  menaces à la validité (\ref{sec:menaces-ref}) et annexe~\ref{ann:incoherences}.
\item \textbf{Chercheur en modélisation} (équipes Inria du programme Apollo, destinataires de la reprise). Jumeau (\ref{sec:jumeau-ref}), sensibilité
  (\ref{sec:sensibilite-ref}), questions ouvertes et module de besoin (\ref{sec:ouvertes-ref}), reproductibilité
  (\ref{sec:repro-ref}).
\item \textbf{Professionnel de la planification.} Ce chapitre, puis données (\ref{sec:donnees-ref}), interface
  (\ref{sec:produit-ref}) et conformité (\ref{sec:conformite-ref}).
\end{itemize}

\subsection{Conventions}
Chaque nombre porte un exposant de provenance :
\begin{itemize}
\item \mesure{} : mesuré sur les données ou par un script du dépôt ;
\item \derive{} : dérivé par calcul de nombres mesurés ou sourcés ;
\item \suppose{} : supposé, c'est-à-dire posé sans mesure ;
\item sourcé : sans exposant, le nombre est suivi d'une citation de la source qui le donne.
\end{itemize}
Les chemins cités sont relatifs à la racine du dépôt. Quand deux fichiers du dépôt donnent des valeurs différentes
pour une même grandeur, le texte retient la mesure la plus récente, et l'annexe~\ref{ann:incoherences} liste les
deux valeurs et la raison du choix. Le rapport décrit et compare ; il ne qualifie pas les territoires (lexique de
\texttt{v2/docs/lab/lexicon.md}).

\subsection{Sigles et codes}\label{int:sigles}
{\small
\begin{description}\setlength{\itemsep}{0pt}
\item[Axe~0] Le jumeau numérique, moteur commun (chapitre~\ref{sec:jumeau-ref}).
\item[Axe~1] Les hyperparamètres sourcés et mesurables (chapitre~\ref{sec:parametres-ref}).
\item[Axe~2] Les méthodes de modélisation et leur banc (chapitre~\ref{sec:methodes-ref}).
\item[Axe~3] La méthodologie de validation et les protocoles pré-enregistrés (chapitre~\ref{sec:validation-ref}).
\item[Axe~4] L'enrichissement des jeux de données (chapitre~\ref{sec:donnees-ref}).
\item[Axe~5] La sensibilité et l'incertitude à 2050 (chapitre~\ref{sec:sensibilite-ref}).
\item[A0] Stock-flux à taux locaux, composition de référence du jumeau (réécriture du modèle de la v1).
\item[A2] Estimateur à rétrécissement hiérarchique de l'axe~2.
\item[A3] À l'axe~2, estimateur à rétention des 60~ans et plus en tendance. À l'axe~5, « méthode A3 de
  l'observatoire » : fourchette $\pm 1{,}6\,a\,h$ reprise d'un outil antérieur de l'équipe.
\item[C1] Estimateur de cohortes bayésiennes de l'axe~2, servi par la carte et l'assistant. Au
  chapitre~\ref{sec:jumeau-ref}, C1 à C18 désignent aussi les critères d'acceptation de \texttt{check.py}.
\item[C2] Variante de C1 qui propage l'incertitude de ses paramètres par approximation de Laplace ; rejetée.
\item[K1] Banc de backtest départemental de la v1 (ancien périmètre libéral + mixte).
\item[N1--N12] Protocoles de validation pré-enregistrés de l'axe~3 (chapitre~\ref{sec:validation-ref}).
\item[DD76] DREES, \emph{Les Dossiers} n\textsuperscript{o}~76 (2021), projections de démographie des
  professions de santé.
\item[ER] \emph{Études et résultats}, publication de la DREES, suivie de son numéro.
\item[RPPS] Répertoire partagé des professionnels de santé.
\item[ADELI] Répertoire antérieur au RPPS ; source des kinésithérapeutes jusqu'en 2016.
\item[SNDS] Système national des données de santé.
\item[BTS] Base Tous Salariés (Insee).
\item[IFSI] Institut de formation en soins infirmiers.
\item[CARPIMKO] Caisse de retraite des infirmiers, kinésithérapeutes, pédicures-podologues, orthophonistes et
  orthoptistes.
\item[COR] Conseil d'orientation des retraites.
\item[MinT] Réconciliation de prévisions hiérarchiques par minimisation de la trace.
\item[CRPS] Score de probabilité classé continu (\emph{continuous ranked probability score}).
\item[SBC] Calibration par simulation (\emph{simulation-based calibration}).
\item[PIT] Transformée intégrale de probabilité (\emph{probability integral transform}).
\item[AIPD] Analyse d'impact relative à la protection des données (RGPD).
\item[PMI] Protection maternelle et infantile.
\item[APL] Accessibilité potentielle localisée (indicateur DREES et IRDES).
\item[2SFCA] Méthode d'accessibilité à deux étapes à aire flottante (\emph{two-step floating catchment area}).
\item[MAPE] Erreur absolue moyenne en pourcentage (\emph{mean absolute percentage error}).
\end{description}
}
